\paragraph{R01. Uniform arithmetic}\label{record-R01}
The common congruence class gives the coefficient cost approximately $1.870614$. The proof includes the mixed 2-adic factorization, odd-prime rank savings and the Chinese remainder construction.

\paragraph{R02. Boundary borders and recurrence lift}\label{record-R02}
The mixed recurrence has two additional boundary moment rows. Coupled first-order integer recurrences generate them from two rational seeds. A larger determinant represents the same forms with scalar cost $O(n\log n)$; its primitive polynomial is unchanged. Two natural borders at $n=24,48$ remain unfavorable. The full matching correction is not rank two at those sizes.

\paragraph{R03. Content of the combined row lattice}\label{record-R03}
The attainable progression is normalized by the rational content of the stacked maximal minors. Its remaining integer step has logarithm $O_x(n\log n)$, so the content determines the leading rate at a fixed rational evaluation.

\paragraph{R04. Prime-scale nonvanishing and covering growth}\label{record-R04}
Three seed determinant polynomials have no rational root, yielding nonvanishing for every fixed rational slope at all sufficiently large prime scales. The same record derives the correction-loss criterion; finite covering growth does not prove it uniformly.

\paragraph{R05. Dual obstruction to bounded approximation}\label{record-R05}
A graph-lattice covering argument certifies the finite approximation for every affine class at the tested sizes, within a prescribed coefficient cap. Exact dual equations show why unrestricted density cannot be assumed in an irrationality proof. The required bounded, uniform covering estimate remains open.

\paragraph{R06. Determinant averaging in the enlarged lattice}\label{record-R06}
Unbiased rounding controls the determinant by a covariance second moment, charging the error directions jointly. This weakens the earlier maximum-norm target. The calculations at $n=48,96$ do not prove the resulting bound as $n$ grows.

\paragraph{R07. Row and column Gram energies}\label{record-R07}
Gram Cauchy--Schwarz separates row and column volumes. A uniform column-energy inequality is proved, and the deterministic endpoint repair is retained in the row Gram. Later records sharpen both quantitative estimates.

\paragraph{R08. Index comparison and a larger coefficient lattice}\label{record-R08}
The uniform index-loss upper bound is $61/288$ on the $n^2$ scale, much larger than the available margin. Removing unnecessary coefficient congruences enlarges each row lattice by an index factor $\exp((5/48+o(1))n^2)$, while retaining the proved arithmetic cost and prime-scale nonvanishing. That extra freedom has no proved uniform real estimate. All 393 tested full-variation-row minors also miss the required rate.

\paragraph{R09. Polynomial kernel inverses}\label{record-R09}
A normalized left-inverse bound controls determinant content but necessarily costs at least the arithmetic threshold. A comparison family has the same dimensions and degree-five profile with a positive quadratic correction loss. Conditional density explains finite successes without establishing irrationality.

\paragraph{R10. Low-degree correction freedom}\label{record-R10}
Corrections below the column band preserve the target period and endpoint zero. A local Smith-image optimization attains the minimum two-layer denominator cost. The first global choices lose other arithmetic and real savings; they are superseded by the uniform class.

\paragraph{R11. Matching the mixed slope}\label{record-R11}
Changing $D$ to $D-16\lambda E$ retains $G+\lambda\pi^2$. Linear high-degree completions satisfying the same band and endpoint conditions give exactly the same determinant. Full, band and endpoint completions were compared with complete primitive normalization.

\paragraph{R12. Moment recurrences and natural bases}\label{record-R12}
The Chebyshev recurrence acquires lower and upper band defects after period matching. Their forced central dyadic defect cannot be removed by the specified integral corrections; a separate rational obstruction rules out a homogeneous extension. Natural recurrence bases, even after removing column content, fail the finite covering criterion.

\paragraph{R13. Moment-derived directions}\label{record-R13}
An interior displacement formula has three boundary terms. Saturated moment-derived directions still fail the finite determinant test. A proposed scalar kernel-pencil route retains a rational approximation error. Degree-five polynomial left inverses of the full low-moment kernel are certified at $n=48,96$; no all-size existence or decay theorem is inferred.

\paragraph{R14. Limit of the absolute-value interpolation bound}\label{record-R14}
Dual probability measures prove that the absolute-value potential method has integrated cost at least 0.04086151228157811175. With the then-current column and repair bounds, it needed to be below 0.037728119494102. This limits that estimate, not the actual row determinant. Sharpening the repair and column bounds later bypasses the limitation.

\paragraph{R15. Three-shift digamma determinants}\label{record-R15}
The determinant and common-node cofactors have no common complex zero for every rational slope. Exact full primitive polynomials still have unfavorable values for the generic tested slopes. Alternative borders, nonzero-sum weights and two four-shift families were exploratory tests, not asymptotic certificates.

\paragraph{R16. The rational image of integer corrections}\label{record-R16}
Smith normal form and an explicit last-row construction show that unrestricted correction determinants attain every value in one arithmetic progression. Mixed-minor content gives its step; the result imposes no height bound on an attaining correction.

\paragraph{R17. Real lattice searches in the paired class}\label{record-R17}
Integer lattice searches within the protected paired-row congruence class give nonzero primitive values below one at $n=24,48,96$. Exact polynomial normalization and interval evaluation certify each example. They do not supply an unbounded irrationality sequence.

\paragraph{R18. The real row operator}\label{record-R18}
The polynomial-part projection is a contraction in the relevant weighted norms. Its failure to commute with the endpoint factor is an explicit finite-rank moment operator. This supplies the main-row plus repair decomposition used in the final estimate.

\paragraph{R19. Repair Gram minors and positive canonical margin}\label{record-R19}
Compound matrices, factorial Schur comparison and a damped energy bound give repair cost below 0.00498. A complete degree-200 column certificate yields a positive canonical margin. The audit covers the stated dependencies, but not later uncontrolled integer corrections.

\paragraph{R20. Spectrum of the main-row operator}\label{record-R20}
Paired rows isolate an explicit Jacobi volume. Rational approximation of the remaining operator bounds its singular values. Twenty-six global potential certificates give spectral cost 0.0450807972 and the main-row rate used in the final estimate.

\paragraph{R21. Paired corrections and protection at two}\label{record-R21}
Copying paired correction rows preserves the original fixed rational rows. Explicit dyadic row divisibility protects the exact determinant valuation at two while allowing most odd-prime savings. The first protected global choices still worsen real decay; later lattice searches address that finite defect.

\paragraph{R22. Structured directions and direct output rounding}\label{record-R22}
Only $b$ independent integer image directions are needed for rounding. The covariance-moment bound for a fixed ordered sequence of image subspaces is optimal over rounding metrics. Twenty-three finite-difference bases at each of $n=48,96$ miss the required margin.

\paragraph{R23. First uniform endpoint-repair bound}\label{record-R23}
Rank-one truncations control every singular value of the repair and give cost at most $0.009486012+0.178211892/n$. This is superseded numerically by the Gram-minor bound, while its kernel identity and projection argument remain in use.

\paragraph{R24. Finite real approximation in the uniform class}\label{record-R24}
The new all-row scalar-modulus class, including divisibility by $2^C$, admits certified small nonzero primitive values at $n=48,96$. These tests use the exact uniform arithmetic family, unlike the earlier paired searches; a uniform real bound remains open.

\paragraph{R25. Individual $\beta(4)$ isolation on an odd grid}\label{record-R25}
Fixed polynomial multipliers cannot isolate $\beta(4)$ in the stated valuation-separated class. Adaptive finite multipliers do isolate it but have unfavorable primitive values. Explicit telescoping multipliers give zero forms, and a finite quotient calculation is kept distinct from an all-degree classification.

\paragraph{R26. Real interpretation of the Padé family}\label{record-R26}
The odd-$m$ Kawashima--Poëls construction gives separate $\beta(4)$ remainders without eliminating $\beta(2)$. A real Cauchy representation, moving half-integer evaluation and complete rational tail conversion give the stated forms and family nonvanishing.

\paragraph{R27. Padé remainder rates and arithmetic cost}\label{record-R27}
A signed saddle calculation gives the actual remainder rate, including eventual nonvanishing. An explicit improved clearing factor still exceeds that decay for every fixed evaluation ratio at least one half. Further primitive coefficient content is not excluded.

\paragraph{R28. Reassessment and changes of target}\label{record-R28}
The useful relaxations are a small set of unresolved even values and rational relations among $1,\pi^2,G$. Simple KP Hankel determinants and lower seed multiplicity were tested with complete coefficients; neither gave favorable growing examples. Other target changes were assessed without claiming new results.

\paragraph{R29. Small-set tails and $\beta(2)$ elimination}\label{record-R29}
Second derivatives with simple zeros can add large-prime constant costs; triple zeros retain grouped cancellation. Undifferentiated forms avoid those primes but their tested $\beta(2)$-elimination weights are costly. All 115 final forms at the tested indices $n\ge3$ have primitive absolute value above one.
